\documentclass[reqno]{amsart} 
\usepackage{amssymb,latexsym,amsmath,amscd,graphicx,setspace,amsthm,verbatim}
\usepackage[margin = 3 cm]{geometry}


\theoremstyle{plain}
\newtheorem{theorem}{Theorem}[section]
\newtheorem{proposition}{Proposition}
\newtheorem{corollary}{Corollary}
\newtheorem{lemma}{Lemma}
\newtheorem{conjecture}{Conjecture}
\newtheorem{question}{Question}
\newtheorem{problem}{Problem}
      
\theoremstyle{definition}
\newtheorem{definition}{Definition}

\newenvironment{solution}{\paragraph{\emph{Solution}.}}{\hfill$\square$}


\newenvironment{solution1}{\paragraph{\emph{Solution $1$}.}}{\hfill$\square$}
\newenvironment{solution2}{\paragraph{\emph{Solution $2$}.}}{\hfill$\square$}
\newenvironment{solution3}{\paragraph{\emph{Solution $3$}.}}{\hfill$\square$}

\begin{document} 

\title[Homework 3]{Homework 3 \\ First due date: September 16.  \\  Due date for the second submission:  September 23.}

\date{\today} 
\maketitle 


\begin{problem}
Consider the function $f:\mathbb{R} \rightarrow \mathbb{R}$ defined via 
\begin{equation*}
x \mapsto f(x) = 
\begin{cases} 
e^{-1/x}, &\text{if } x>0;\\
0, &\text{if } x \le 0.
\end{cases} 
\end{equation*}
\begin{enumerate}
\item Show that $f$ is differentiable everywhere.
\item What is $f'(0)$?
\item Using an induction argument, show that $f^{(n)}(0) = 0$ for all $n \ge 0$.
\item Is $f$ smooth at $x=0$?  Explain carefully.
\item Is $f$ analytic at $x=0$?  Explain carefully.
\end{enumerate}
\end{problem}
\begin{solution}

\end{solution}

In class, we defined a function $f:\mathbb{R}^{m} \rightarrow \mathbb{R}^{n}$ to be differentiable at $p \in \mathbb{R}^{m}$ if there exists a linear transformation $L:\mathbb{V}^{m} \rightarrow \mathbb{V}^{n}$ such that for all $\vec{v} \in \mathbb{V}^{m}$, one has
$$f(p + \vec{v}) - f(p) = L(\vec{v}) + \varepsilon(\vec{v}), $$
where $\varepsilon:\mathbb{V}^{m} \rightarrow \mathbb{V}^{n}$ is a function (an error term) satisfying
$$\frac{||\varepsilon(\vec{v})||}{||\vec{v}||} \to 0 $$
as $\vec{v} \to \vec{0}$.  Moreover, we showed that this linear transformation $L$ is unique if it exists.  It is called the derivative of $f$ at $p$ and is denoted by $Df_{p}$ or a similar notation.  In the next problem, you will go through such an example in details.

\begin{problem}
Throughout this problem, we let $f:\mathbb{R}^{2} \rightarrow \mathbb{R}^{2}$ be defined via
$$(x,y) \mapsto f(x,y) = (x,x^{2} + y^{2}). $$
\begin{enumerate}
\item Consider the function $L:\mathbb{V}^{2} \rightarrow \mathbb{V}^{2}$ defined via
$$\vec{v} = \langle v_{1},v_{2} \rangle \mapsto L(\vec{v}) = \langle v_{1},2v_{1} + v_{2} \rangle. $$
Put your Math 235 hat on and show that $L$ is a linear transformation.
\item Let now $p = (1,1)$.  For an arbitrary $\vec{v} = \langle v_{1}, v_{2} \rangle$ in $\mathbb{V}^{2}$, calculate
$$f(p + \vec{v}) - f(p) - L(\vec{v}) \in \mathbb{V}^{2}. $$
(You should get a two dimensional vector with coordinates expressed in terms of $v_{1}$ and $v_{2}$.)
\item Let us denote the last expression $f(p + \vec{v}) - f(p) - L(\vec{v})$ by $\varepsilon(\vec{v})$.  (Note that $\varepsilon$ is a function from $\mathbb{V}^{2}$ to $\mathbb{V}^{2}$.)  Prove carefully that
$$\frac{||\varepsilon(\vec{v})||}{||\vec{v}||} \to 0 $$
as $\vec{v} \to \vec{0}$.
\item Is $f$ differentiable at $p = (1,1)$?  Explain carefully.
\item What is the derivative of $f$ at $p = (1,1)$?
\end{enumerate}
\end{problem}
\begin{solution}

\end{solution}

The next problem will be useful for a topic we will discuss in class this coming week.  If you don't remember what it means for a function $f:A \rightarrow B$ to be \emph{injective} from Math 330, look it up first.

\begin{problem}
Consider the function $\gamma:\left(\frac{\pi}{2},\frac{5 \pi}{2} \right) \rightarrow \mathbb{R}^{2}$ defined via
$$\theta \mapsto \gamma(\theta) = \left(\frac{\cos(\theta)}{1+\sin^{2}(\theta)},\frac{\sin(\theta) \cos(\theta)}{1+\sin^{2}(\theta)} \right). $$
Show carefully that $\gamma$ is injective.
\end{problem}
\begin{solution}

\end{solution}

\flushleft
{\bfseries Problems for graduate students (bonus points for undergraduate students)}

\begin{problem}
Show that a function $f:\mathbb{R}^{m} \rightarrow \mathbb{R}^{n}$ is continuous at $p \in \mathbb{R}^{m}$ if and only if all coordinate functions $f_{i}:\mathbb{R}^{m} \rightarrow \mathbb{R}$ for $i=1,\ldots,n$ are continuous at $p \in \mathbb{R}^{m}$.  (Recall:  a function $f:\mathbb{R}^{m} \rightarrow \mathbb{R}^{n}$ is continuous at $p \in \mathbb{R}^{m}$ if for all $\varepsilon > 0$, there exists $\delta > 0$ such that if $d_{\mathbb{R}^{m}}(x,p)< \delta$, then $d_{\mathbb{R}^{n}}(f(x),f(p))< \varepsilon$.)
\end{problem}
\begin{solution}

\end{solution}

\begin{problem}
Show carefully that if $f:\mathbb{R}^{m} \rightarrow \mathbb{R}^{n}$ is differentiable at $p \in \mathbb{R}^{m}$ (according to the definition above), then $f$ is continuous at $p \in \mathbb{R}^{m}$.
\end{problem}
\begin{solution}

\end{solution}

\begin{problem}
Show carefully that a function $f:\mathbb{R}^{m} \rightarrow \mathbb{R}^{n}$ is differentiable at $p \in \mathbb{R}^{m}$ if and only if all coordinate functions $f_{i}:\mathbb{R}^{m} \rightarrow \mathbb{R}$ for $i=1,\ldots,n$ are differentiable at $p \in \mathbb{R}^{m}$.  (Again: use the definition of differentiability from above.)   
\end{problem}
\begin{solution}

\end{solution}











\end{document} 



